\documentclass[UTF8,12pt,a4paper]{ctexart}

\usepackage{amsmath,amssymb}
\usepackage{geometry}
\usepackage{booktabs}
\usepackage{tabularx}
\usepackage{array}
\usepackage{enumitem}
\usepackage{graphicx}
\usepackage{float}
\usepackage{caption}
\usepackage{fancyhdr}
\usepackage{hyperref}
\usepackage{pgfplots}
\usepackage{fontspec}

\geometry{left=2.5cm,right=2.5cm,top=2.5cm,bottom=2.5cm}
\IfFontExistsTF{Times New Roman}{\setmainfont{Times New Roman}}{}
\IfFontExistsTF{SimSun}{
    \IfFontExistsTF{SimHei}{\setCJKmainfont[BoldFont=SimHei]{SimSun}}{\setCJKmainfont{SimSun}}
}{}
\pgfplotsset{compat=1.18}
\hypersetup{hidelinks}
\captionsetup{font=small,labelsep=quad}

\pagestyle{fancy}
\fancyhf{}
\lhead{RLC 电路实验}
\rhead{基础物理实验}
\cfoot{\thepage}
\setlength{\headheight}{14.5pt}
\renewcommand{\headrulewidth}{0.4pt}

\setlength{\parindent}{2em}
\setlength{\parskip}{0.2em}
\linespread{1.2}
\numberwithin{equation}{section}

\newcolumntype{Y}{>{\raggedright\arraybackslash}X}
\newcolumntype{C}{>{\centering\arraybackslash}X}

\begin{document}

\begin{center}
    {\zihao{2}\bfseries RLC 电路实验报告}
\end{center}

\vspace{0.5em}

\begin{center}
\begin{tabular}{rlrl}
实验名称： & RLC 电路实验 & 课程： & 基础物理实验 \\
姓名： & 纪文龙 & 班级： & 行健烽火 4 \\
学号： & 2024012842 & 实验日期： & 2025 年 3 月 9 日 \\
\end{tabular}
\end{center}

\section{实验目的}

\begin{enumerate}[label=\arabic*.,leftmargin=2.2em]
    \item 观察 RLC 串联电路在方波激励下的阻尼振荡现象，区分欠阻尼、临界阻尼和过阻尼三种状态。
    \item 测量欠阻尼振荡的衰减时间常数和阻尼系数，并与由回路等效电阻得到的理论阻尼系数比较。
    \item 研究 RLC 串联电路在正弦信号驱动下的受迫振荡特性，测量幅频响应曲线和相频响应曲线，确定共振频率、半功率带宽和品质因数。
    \item 在共振频率附近核算信号源输出功率与电阻性元件消耗功率的关系，判断功率账自洽情况。
\end{enumerate}

\section{实验原理}

RLC 串联电路是典型的二阶线性电路。以电容两端电压 \(v_C(t)\) 为变量，在无外加驱动时由基尔霍夫电压定律可得
\begin{equation}
    \frac{\mathrm{d}^2 v_C}{\mathrm{d}t^2}
    +\frac{R_\Sigma}{L}\frac{\mathrm{d}v_C}{\mathrm{d}t}
    +\frac{1}{LC}v_C=0,
\end{equation}
其中 \(R_\Sigma\) 表示该状态下回路中的等效总电阻，包含外接电阻和相关损耗电阻。定义
\begin{equation}
    \beta=\frac{R_\Sigma}{2L},\qquad
    \omega_0=\frac{1}{\sqrt{LC}}.
\end{equation}

当 \(\beta<\omega_0\) 时，电路处于欠阻尼状态，电容电压可写为
\begin{equation}
    v_C(t)=A e^{-\beta t}\cos(\omega t+\varphi_0),
    \qquad
    \omega=\sqrt{\omega_0^2-\beta^2}.
\end{equation}
其振幅包络满足
\begin{equation}
    V_C(t)=A e^{-\beta t},\qquad
    \ln V_C=\ln A-\beta t.
\end{equation}
因此，对逐次峰谷电压的绝对值取对数并作线性拟合，可由斜率求得实验阻尼系数。临界阻尼条件为 \(\beta=\omega_0\)，对应临界电阻
\begin{equation}
    R_c=2\sqrt{\frac{L}{C}}.
\end{equation}

当电路受到正弦电压
\[
    v_i(t)=V_m\cos\omega t
\]
驱动时，稳态电流峰值和电流相对于信号源电压的相位差为
\begin{equation}
    I_m=\frac{V_m}{\sqrt{R_\Sigma^2+\left(\omega L-\frac{1}{\omega C}\right)^2}},
\end{equation}
\begin{equation}
    \varphi=-\arctan\frac{\omega L-\frac{1}{\omega C}}{R_\Sigma}.
\end{equation}
当 \(\omega\approx\omega_0\) 时，感抗和容抗近似抵消，电流达到最大值，相位差接近 \(0^\circ\)，电路表现为近似纯电阻性。理论共振频率为
\begin{equation}
    f_{0,\mathrm{th}}=\frac{1}{2\pi\sqrt{LC}}.
\end{equation}
品质因数可由半功率带宽计算：
\begin{equation}
    \Delta f=f_+-f_-,
    \qquad
    Q=\frac{f_0}{\Delta f}.
\end{equation}
若使用电压、电流峰值，则平均功率统一写为
\begin{equation}
    P=\frac{V_m I_m\cos\varphi}{2}.
\end{equation}

\section{实验仪器与主要参数}

实验仪器包括示波器、信号发生器、数字万用表、\(5\ \mathrm{k}\Omega\) 可变电阻、\(20\ \mathrm{mH}\) 电感、\(0.01\ \mu\mathrm{F}\) 电容及连接导线。

本文采用的主要元件参数为
\begin{equation}
    L=20\ \mathrm{mH}=0.020\ \mathrm{H},\qquad
    C=0.01\ \mu\mathrm{F}=1.0\times10^{-8}\ \mathrm{F}.
\end{equation}
原始记录中电感直流内阻为
\begin{equation}
    R_L=92.825\ \Omega.
\end{equation}
该值反映电感直流电阻，可用于误差讨论，但不直接等同于受迫振荡分压估算得到的等效损耗。

\section{实验任务及步骤}

\begin{enumerate}[label=\arabic*.,leftmargin=2.2em]
    \item 按实验线路连接 RLC 串联电路，用方波信号驱动电路，调节可变电阻，观察并记录欠阻尼、临界阻尼和过阻尼三种状态下 \(v_C(t)\) 的波形。
    \item 调至临界阻尼附近，记录临界阻尼对应的电阻读数；若原始记录未保留该读数，则只给出理论临界电阻作为参考。
    \item 将电路调至欠阻尼状态，记录逐次峰谷电压 \(V_C\) 及相关时间读数，对 \(\ln(V_C)\) 与测量序号或时间进行线性拟合，求阻尼系数 \(\beta\)。
    \item 以正弦波驱动 RLC 串联电路，在 \(5\sim15\ \mathrm{kHz}\) 范围内扫频，记录电阻电压峰值 \(V_{Rm}\) 和相位差 \(\varphi\)，绘制幅频响应曲线和相频响应曲线，确定 \(f_0\)、\(\Delta f\) 与 \(Q\)。
    \item 在共振频率附近测量各元件电压及相位，统一使用峰值功率公式，核算信号源输出功率、电阻功率和损耗功率之间的关系。
\end{enumerate}

\section{数据处理}

\subsection{符号与取值说明}

为避免一个 \(R\) 在不同环节中含义跳变，本文采用表~\ref{tab:symbols} 所示符号约定。

\begin{table}[H]
\centering
\caption{主要符号与取值说明}
\label{tab:symbols}
\small
\begin{tabularx}{\textwidth}{>{\centering\arraybackslash}p{2.5cm} Y Y >{\raggedleft\arraybackslash}p{2.8cm}}
\toprule
符号 & 物理含义与使用位置 & 说明 & 取值 \\
\midrule
\(L\) & 电感标称值，全文理论计算 & 来自实验仪器 & \(20\ \mathrm{mH}\) \\
\(C\) & 电容标称值，全文理论计算 & 来自实验仪器 & \(0.01\ \mu\mathrm{F}\) \\
\(R_L\) & 电感直流内阻，误差讨论 & 不直接等同于交流等效损耗 & \(92.825\ \Omega\) \\
\(R_d\) & 阻尼振荡外接电阻 & 用于 \(150\ \Omega+\) 寄生电阻链条 & \(150\ \Omega\) \\
\(R_f\) & 频响扫描图题标称电阻 & 原始频响表和图题标注值 & \(2.5\ \mathrm{k}\Omega\) \\
\(R_m\) & 分压估算和电流计算所用测量电阻 & 按教师在分压公式旁批注采用 & \(500\ \Omega\) \\
\(R_{\mathrm{loss},d}\) & 阻尼振荡原始损耗 & 用于主写 \(\beta_{\mathrm{th}}\) & \(140\ \Omega\) \\
\(R_{\mathrm{loss},f}\) & 受迫振荡分压估算损耗 & 由 \(R_m=500\ \Omega\)、\(V_m/V_R\) 反推 & \(194.3\ \Omega\) \\
\bottomrule
\end{tabularx}
\end{table}

需要强调的是，\(R_{\mathrm{loss},d}=140\ \Omega\) 与 \(R_{\mathrm{loss},f}=194.3\ \Omega\) 来自不同数据链。前者用于阻尼振荡理论阻尼系数比较；后者只用于受迫振荡分压估算和共振处功率账自洽检查。本文不将二者合并为单一损耗电阻。

\subsection{电源操作}

电源设定为恒流 \(I=0.1\ \mathrm{A}\)、限压 \(V=5\ \mathrm{V}\)。改变负载电阻得到表~\ref{tab:power} 所示数据。

\begin{table}[H]
\centering
\caption{电源操作测试数据}
\label{tab:power}
\begin{tabular}{cccc}
\toprule
电阻 \(R/\Omega\) & 电压 \(V/\mathrm{V}\) & 电流 \(I/\mathrm{A}\) & 工作状态 \\
\midrule
30 & 3.19 & 0.100 & CC，恒流 \\
40 & 4.18 & 0.100 & CC，恒流 \\
50 & 5.00 & 0.095 & CV，恒压 \\
60 & 5.00 & 0.081 & CV，恒压 \\
70 & 5.00 & 0.070 & CV，恒压 \\
\bottomrule
\end{tabular}
\end{table}

当 \(R=30\ \Omega\) 和 \(40\ \Omega\) 时，维持 \(0.1\ \mathrm{A}\) 所需电压分别约为 \(3.0\ \mathrm{V}\) 和 \(4.0\ \mathrm{V}\)，低于限压 \(5\ \mathrm{V}\)，电源处于恒流状态。当 \(R\ge 50\ \Omega\) 时，维持 \(0.1\ \mathrm{A}\) 所需电压达到或超过限压，电源转入恒压状态，电流随负载电阻增大而下降。

\subsection{阻尼振荡}

\subsubsection{临界阻尼电阻}

由 \(L=0.020\ \mathrm{H}\)、\(C=1.0\times10^{-8}\ \mathrm{F}\) 得
\begin{align}
    R_c
    &=2\sqrt{\frac{L}{C}}
      =2\sqrt{\frac{0.020}{1.0\times10^{-8}}} \notag\\
    &=2.828\times10^3\ \Omega
      \approx 2.83\ \mathrm{k}\Omega.
\end{align}
原始数据记录中没有保留临界阻尼电阻的实测值，因此此处只给出理论值作为波形判断参考。

\subsubsection{欠阻尼峰谷数据与对数修正}

欠阻尼振荡中，原始记录给出第 1 次和第 6 次峰谷时间读数
\[
    t_1=76\ \mu\mathrm{s},\qquad t_6=502\ \mu\mathrm{s}.
\]
逐次记录的相邻峰谷间隔按一个周期 \(T\) 处理，因此
\begin{equation}
    T=\frac{t_6-t_1}{6-1}
    =\frac{502-76}{5}\ \mu\mathrm{s}
    =85.2\ \mu\mathrm{s}.
\end{equation}

表~\ref{tab:damping-data} 中第 6 个点采用由 \(V_C=0.0781\ \mathrm{V}\) 重新计算得到的对数值：
\[
    \ln(0.0781)=-2.5498.
\]

\begin{table}[H]
\centering
\caption{欠阻尼振荡峰谷电压及对数值}
\label{tab:damping-data}
\begin{tabular}{cccc}
\toprule
测量次数 \(n\) & \(V_C/\mathrm{V}\) & \(\ln(V_C)\) & 时间记录 \\
\midrule
1 & 4.7180 & 1.5514 & \(76\ \mu\mathrm{s}\) \\
2 & 1.9532 & 0.6695 & -- \\
3 & 0.8437 & \(-0.1700\) & -- \\
4 & 0.3906 & \(-0.9401\) & -- \\
5 & 0.1781 & \(-1.7254\) & -- \\
6 & 0.0781 & \(-2.5498\) & \(502\ \mu\mathrm{s}\) \\
\bottomrule
\end{tabular}
\end{table}

对 \(\ln(V_C)\) 与测量次数 \(n\) 作最小二乘线性拟合，可得
\begin{equation}
    \ln(V_C)=-0.8132n+2.3187,\qquad r=-0.9998.
\end{equation}
数据点与拟合直线如图~\ref{fig:damping-fit} 所示。

\begin{figure}[H]
\centering
\begin{tikzpicture}
\begin{axis}[
    width=0.82\textwidth,
    height=6cm,
    grid=both,
    xlabel={测量次数 \(n\)},
    ylabel={\(\ln(V_C)\)},
    xmin=0.6,xmax=6.4,
    ymin=-3.0,ymax=2.0,
    legend pos=south west
]
\addplot+[only marks,mark=*,mark size=2.2pt] coordinates {
    (1,1.5514) (2,0.6695) (3,-0.1700)
    (4,-0.9401) (5,-1.7254) (6,-2.5498)
};
\addlegendentry{实验数据}
\addplot+[domain=1:6,samples=2,thick] {-0.8132*x+2.3187};
\addlegendentry{线性拟合}
\end{axis}
\end{tikzpicture}
\caption{\(\ln(V_C)\) 与测量次数 \(n\) 的线性拟合}
\label{fig:damping-fit}
\end{figure}

欠阻尼振荡的包络满足 \(\ln V_C=\ln A-\beta t\)。由于 \(n\) 每增加 1 对应时间增加 \(T=85.2\ \mu\mathrm{s}\)，拟合斜率的绝对值等于 \(\beta T\)，于是
\begin{equation}
    \beta_{\mathrm{fit}}
    =\frac{0.8132}{85.2\times10^{-6}}
    =9.54\times10^3\ \mathrm{s^{-1}}.
\end{equation}
对应衰减时间常数为
\begin{equation}
    \tau=\frac{1}{\beta_{\mathrm{fit}}}
    =1.048\times10^{-4}\ \mathrm{s}
    =104.8\ \mu\mathrm{s}.
\end{equation}

\subsubsection{阻尼系数理论值比较}

阻尼振荡理论阻尼系数应使用阻尼振荡回路中的等效总电阻，而不是受迫振荡频响图题中的 \(2.5\ \mathrm{k}\Omega\)。根据原始记录和教师批注，取
\[
    R_d=150\ \Omega,\qquad
    R_{\mathrm{loss},d}=140\ \Omega.
\]
则
\begin{equation}
    \beta_{\mathrm{th}}
    =\frac{R_d+R_{\mathrm{loss},d}}{2L}
    =\frac{150+140}{2\times0.020}
    =7.25\times10^3\ \mathrm{s^{-1}}.
\end{equation}
实验拟合值与理论值的相对差异为
\begin{equation}
    \frac{\beta_{\mathrm{fit}}-\beta_{\mathrm{th}}}{\beta_{\mathrm{th}}}\times100\%
    =\frac{9.54-7.25}{7.25}\times100\%
    =31.6\%.
\end{equation}
拟合得到的衰减快于由 \(150\ \Omega+140\ \Omega\) 估计的理论衰减。差异可能来自信号源内阻、接触电阻、电感交流损耗和示波器读数误差等影响。

若把受迫振荡分压估算值 \(R_{\mathrm{loss},f}=194.3\ \Omega\) 条件性代入阻尼公式，可得
\begin{equation}
    \beta_{\mathrm{th}}'
    =\frac{150+194.3}{2\times0.020}
    =8.61\times10^3\ \mathrm{s^{-1}},
\end{equation}
此时与拟合值的差异约为 \(10.9\%\)。该结果只作为条件性对照，不作为阻尼振荡理论比较的主结果。

\subsection{受迫振荡}

\subsubsection{原始频响数据与共振频率}

受迫振荡中，原始频响表和图题标注的扫描电阻为
\[
    R_f=2.5\ \mathrm{k}\Omega.
\]
输入电压峰值记录为
\[
    V_m=4.621\ \mathrm{V}.
\]
频率响应数据如表~\ref{tab:freq-data} 所示。

\begin{table}[H]
\centering
\caption{受迫振荡频率响应数据}
\label{tab:freq-data}
\begin{tabular}{ccc}
\toprule
频率 \(f/\mathrm{kHz}\) & \(V_{Rm}/\mathrm{V}\) & 相位差 \(\varphi/^\circ\) \\
\midrule
5.0 & 0.933 & \(-73.5\) \\
6.0 & 1.211 & \(-69.9\) \\
7.0 & 1.600 & \(-64.4\) \\
8.0 & 2.029 & \(-55.8\) \\
9.0 & 2.641 & \(-43.2\) \\
10.0 & 3.017 & \(-26.1\) \\
10.5 & 3.187 & \(-11.6\) \\
11.0 & 3.287 & \(-3.7\) \\
11.5 & 3.311 & 4.0 \\
12.0 & 3.268 & 9.4 \\
13.0 & 3.055 & 22.9 \\
14.0 & 2.770 & 17.2 \\
15.0 & 2.502 & 23.6 \\
\bottomrule
\end{tabular}
\end{table}

由表中离散数据可见，电阻电压峰值在 \(11.5\ \mathrm{kHz}\) 附近达到表内最大值 \(3.311\ \mathrm{V}\)。原图和正文标注的共振频率为 \(11.4\ \mathrm{kHz}\)，可解释为由响应曲线读数或插值得到。相频曲线在 \(11.0\sim11.5\ \mathrm{kHz}\) 之间穿过 \(0^\circ\)，也支持共振频率位于该区间。

\begin{figure}[H]
\centering
\begin{tikzpicture}
\begin{axis}[
    width=0.82\textwidth,
    height=6cm,
    grid=both,
    xlabel={频率 \(f/\mathrm{kHz}\)},
    ylabel={\(V_{Rm}/\mathrm{V}\)},
    xmin=4.8,xmax=15.6,
    ymin=0.5,ymax=3.6,
    legend pos=south east
]
\addplot+[mark=*,mark size=1.8pt,thick] coordinates {
    (5.0,0.933) (6.0,1.211) (7.0,1.600) (8.0,2.029)
    (9.0,2.641) (10.0,3.017) (10.5,3.187) (11.0,3.287)
    (11.5,3.311) (12.0,3.268) (13.0,3.055) (14.0,2.770) (15.0,2.502)
};
\addlegendentry{\(V_{Rm}\)}
\addplot+[domain=5:15.5,samples=2,dashed] {2.34};
\addlegendentry{半功率电压}
\addplot+[domain=0.5:3.6,samples=2,densely dotted] coordinates {(11.4,0.5) (11.4,3.6)};
\addlegendentry{\(f_0=11.4\ \mathrm{kHz}\)}
\end{axis}
\end{tikzpicture}
\caption{受迫振荡幅频响应曲线}
\label{fig:vrm-freq}
\end{figure}

\begin{figure}[H]
\centering
\begin{tikzpicture}
\begin{axis}[
    width=0.82\textwidth,
    height=6cm,
    grid=both,
    xlabel={频率 \(f/\mathrm{kHz}\)},
    ylabel={相位差 \(\varphi/^\circ\)},
    xmin=4.8,xmax=15.6,
    ymin=-80,ymax=30,
    legend pos=south east
]
\addplot+[mark=*,mark size=1.8pt,thick] coordinates {
    (5.0,-73.5) (6.0,-69.9) (7.0,-64.4) (8.0,-55.8)
    (9.0,-43.2) (10.0,-26.1) (10.5,-11.6) (11.0,-3.7)
    (11.5,4.0) (12.0,9.4) (13.0,22.9) (14.0,17.2) (15.0,23.6)
};
\addlegendentry{\(\varphi\)}
\addplot+[domain=5:15.5,samples=2,dashed] {0};
\addlegendentry{\(\varphi=0^\circ\)}
\addplot+[domain=-80:30,samples=2,densely dotted] coordinates {(11.4,-80) (11.4,30)};
\addlegendentry{\(f_0=11.4\ \mathrm{kHz}\)}
\end{axis}
\end{tikzpicture}
\caption{受迫振荡相频响应曲线}
\label{fig:phi-freq}
\end{figure}

理论共振频率为
\begin{align}
    f_{0,\mathrm{th}}
    &=\frac{1}{2\pi\sqrt{LC}}
      =\frac{1}{2\pi\sqrt{0.020\times1.0\times10^{-8}}} \notag\\
    &=1.1254\times10^4\ \mathrm{Hz}
      =11.254\ \mathrm{kHz}.
\end{align}
采用曲线读数
\[
    f_{0,\mathrm{exp}}=11.4\ \mathrm{kHz},
\]
则相对误差为
\begin{equation}
    \delta
    =\frac{|11.4-11.254|}{11.254}\times100\%
    =1.30\%.
\end{equation}

\subsubsection{半功率带宽与品质因数}

由幅频曲线最大值 \(V_{Rm,\max}\approx3.31\ \mathrm{V}\) 可得半功率电压参考值
\begin{equation}
    \frac{V_{Rm,\max}}{\sqrt{2}}
    \approx \frac{3.311}{\sqrt{2}}
    =2.34\ \mathrm{V}.
\end{equation}
原始记录给出的半功率频率为
\[
    f_-=8.7\ \mathrm{kHz},\qquad
    f_+=15.4\ \mathrm{kHz}.
\]
其中 \(f_+=15.4\ \mathrm{kHz}\) 未直接出现在频响表中，应理解为由曲线读数、表外原始记录或外推得到。由此
\begin{equation}
    \Delta f=f_+-f_-
    =15.4-8.7
    =6.7\ \mathrm{kHz},
\end{equation}
\begin{equation}
    Q_{\mathrm{exp}}
    =\frac{f_{0,\mathrm{exp}}}{\Delta f}
    =\frac{11.4}{6.7}
    =1.70.
\end{equation}

\subsubsection{图题电阻与批注电阻的条件性处理}

原始频响表及图题标注的扫描电阻为
\[
    R_f=2.5\ \mathrm{k}\Omega.
\]
但在寄生电阻分压公式处，教师明确批注该式中的测量电阻应取
\[
    R_m=500\ \Omega.
\]
因此本文将图题中的标称扫描电阻记为 \(R_f\)，将分压估算与共振点电流计算所用电阻单独记为 \(R_m=500\ \Omega\)。鉴于两处原始记录存在冲突，以下不将二者混写为同一实验设置下的同一个 \(R\)，相关结论按条件性结果处理。

这一冲突也可由带宽作数量级检查。由实验带宽反推串联总电阻，有
\begin{equation}
    R_{\Delta f}\approx 2\pi L\Delta f
    =2\pi\times0.020\times6.7\times10^3
    \approx842\ \Omega.
\end{equation}
该数量级明显小于 \(2.5\ \mathrm{k}\Omega\)，而与 \(R_m+R_{\mathrm{loss},f}\) 的量级更接近。因此，本文保留图题中的标称扫描电阻 \(R_f=2.5\ \mathrm{k}\Omega\)，同时将分压估算和功率核算中的测量电阻按教师批注单独记为 \(R_m=500\ \Omega\)。

\subsubsection{受迫振荡分压估算损耗}

共振附近记录
\[
    V_m=4.621\ \mathrm{V},\qquad
    V_R=3.328\ \mathrm{V}.
\]
按教师批注，分压公式中的测量电阻取 \(R_m=500\ \Omega\)，则受迫振荡分压估算损耗为
\begin{align}
    R_{\mathrm{loss},f}
    &=R_m\left(\frac{V_m}{V_R}-1\right) \notag\\
    &=500\left(\frac{4.621}{3.328}-1\right)
      =194.3\ \Omega.
\end{align}
若误用图题中的 \(2.5\ \mathrm{k}\Omega\) 代入同一公式，会得到约 \(971\ \Omega\)。本文将 \(194.3\ \Omega\) 仅记为受迫振荡分压估算损耗 \(R_{\mathrm{loss},f}\)，不把它改写成全文唯一的寄生电阻。

\subsection{功率核算}

共振频率取
\[
    f_0=11.4\ \mathrm{kHz},\qquad
    \omega_0=2\pi f_0=2\pi\times11.4\times10^3.
\]
本节统一采用峰值体系。共振点电阻电压为 \(V_R=3.328\ \mathrm{V}\)，按教师批注取 \(R_m=500\ \Omega\)，电流峰值为
\begin{equation}
    I_m=\frac{V_R}{R_m}
    =\frac{3.328}{500}
    =6.656\times10^{-3}\ \mathrm{A}
    =6.656\ \mathrm{mA}.
\end{equation}
由此得到电容和理想电感上的电压峰值
\begin{align}
    V_C
    &=\frac{I_m}{\omega_0 C}
      =\frac{6.656\times10^{-3}}
      {2\pi\times11.4\times10^3\times1.0\times10^{-8}}
      =9.29\ \mathrm{V},\\
    V_L
    &=I_m\omega_0 L
      =6.656\times10^{-3}\times2\pi\times11.4\times10^3\times0.020
      =9.54\ \mathrm{V}.
\end{align}
两者幅值接近、相位相反，说明在共振附近理想电感与电容主要交换能量，平均功率近似为零。

使用峰值功率公式，功率核算如表~\ref{tab:power-account} 所示。

\begin{table}[H]
\centering
\caption{共振点功率核算}
\label{tab:power-account}
\begin{tabularx}{\textwidth}{Y C C}
\toprule
对象 & 计算式 & 结果 \\
\midrule
电流峰值 & \(I_m=3.328/500\) & \(6.656\ \mathrm{mA}\) \\
外接测量电阻功率 & \(P_R=V_R I_m/2\) & \(11.08\ \mathrm{mW}\) \\
受迫估算损耗功率 & \(P_{\mathrm{loss},f}=I_m^2R_{\mathrm{loss},f}/2\) & \(4.30\ \mathrm{mW}\) \\
信号源输出功率 & \(P_{\mathrm{source}}=V_m I_m/2\) & \(15.38\ \mathrm{mW}\) \\
功率账 & \(P_R+P_{\mathrm{loss},f}\) & \(15.38\ \mathrm{mW}\) \\
\bottomrule
\end{tabularx}
\end{table}

在采用同一组 \(V_m,V_R,R_m=500\ \Omega\) 反推 \(R_{\mathrm{loss},f}=194.3\ \Omega\) 后，有
\begin{equation}
    P_R+P_{\mathrm{loss},f}\approx P_{\mathrm{source}}.
\end{equation}
但这个等式只能说明修正后的功率账在代数上自洽，不能作为独立检验。原因是 \(R_{\mathrm{loss},f}\) 本身就是由 \(V_m/V_R\) 反推得到，代回功率式会自然闭合。

若使用独立记录的阻尼损耗 \(R_{\mathrm{loss},d}=140\ \Omega\) 作对照，则
\begin{equation}
    P_{\mathrm{loss},d}
    =\frac{I_m^2R_{\mathrm{loss},d}}{2}
    =3.10\ \mathrm{mW},
\end{equation}
\begin{equation}
    P_R+P_{\mathrm{loss},d}
    =11.08+3.10
    =14.18\ \mathrm{mW}.
\end{equation}
与信号源输出功率 \(P_{\mathrm{source}}=15.38\ \mathrm{mW}\) 相比，偏差为
\begin{equation}
    \frac{|15.38-14.18|}{15.38}\times100\%
    =7.8\%.
\end{equation}
因此，本节结论限于：采用分压估算损耗时功率账自洽；独立损耗对照显示功率账存在约 \(7.8\%\) 的差异，差异主要来自损耗电阻等效、信号源内阻、相位读数和元件交流损耗。

\subsection{核心结果汇总}

\begin{table}[H]
\centering
\caption{核心结果汇总}
\label{tab:summary}
\small
\begin{tabularx}{\textwidth}{>{\centering\arraybackslash}p{2.2cm} >{\centering\arraybackslash}p{3.2cm} >{\raggedleft\arraybackslash}p{3.2cm} Y}
\toprule
模块 & 关键量 & 结果 & 说明 \\
\midrule
阻尼振荡 & \(T\) & \(85.2\ \mu\mathrm{s}\) & 相邻记录点按一个周期处理 \\
阻尼振荡 & \(\beta_{\mathrm{fit}}\) & \(9.54\times10^3\ \mathrm{s^{-1}}\) & 由 \(\ln(V_C)=-0.8132n+2.3187\) 得到 \\
阻尼振荡 & \(\tau\) & \(104.8\ \mu\mathrm{s}\) & \(\tau=1/\beta\) \\
阻尼振荡 & \(\beta_{\mathrm{th}}\) & \(7.25\times10^3\ \mathrm{s^{-1}}\) & 主写，取 \(R_d+R_{\mathrm{loss},d}=150+140\ \Omega\) \\
阻尼振荡 & \(\beta_{\mathrm{th}}'\) & \(8.61\times10^3\ \mathrm{s^{-1}}\) & 仅作条件性对照 \\
受迫振荡 & \(f_{0,\mathrm{th}}\) & \(11.254\ \mathrm{kHz}\) & 由 \(1/(2\pi\sqrt{LC})\) 得到 \\
受迫振荡 & \(f_{0,\mathrm{exp}}\) & \(11.4\ \mathrm{kHz}\) & 曲线读数或插值，误差 \(1.30\%\) \\
受迫振荡 & \(\Delta f\) & \(6.7\ \mathrm{kHz}\) & \(f_-=8.7\ \mathrm{kHz}\)，\(f_+=15.4\ \mathrm{kHz}\) \\
受迫振荡 & \(Q\) & \(1.70\) & \(Q=f_0/\Delta f\) \\
损耗电阻 & \(R_{\mathrm{loss},d}\) & \(140\ \Omega\) & 阻尼振荡原始损耗 \\
损耗电阻 & \(R_{\mathrm{loss},f}\) & \(194.3\ \Omega\) & 受迫振荡分压估算损耗 \\
功率核算 & \(I_m\) & \(6.656\ \mathrm{mA}\) & 按 \(R_m=500\ \Omega\) \\
功率核算 & \(P_R\) & \(11.08\ \mathrm{mW}\) & 外接测量电阻功率 \\
功率核算 & \(P_{\mathrm{loss},f}\) & \(4.30\ \mathrm{mW}\) & 受迫估算损耗功率 \\
功率核算 & \(P_{\mathrm{source}}\) & \(15.38\ \mathrm{mW}\) & 与 \(P_R+P_{\mathrm{loss},f}\) 构成功率账自洽 \\
\bottomrule
\end{tabularx}
\end{table}

\section{误差分析}

\begin{enumerate}[label=\arabic*.,leftmargin=2.2em]
    \item 电阻取值定义不清会直接影响电流、损耗和功率。原始材料中同时出现 \(R_f=2.5\ \mathrm{k}\Omega\) 和教师批注 \(R_m=500\ \Omega\)。若在 \(I_m=V_R/R_m\) 中误用 \(2.5\ \mathrm{k}\Omega\)，电流会被低估为正确值的约五分之一，电容、电感电压和平均功率也会随之偏小。
    \item 等效损耗电阻会影响 \(\beta\)、\(Q\) 和带宽。对阻尼振荡，等效总电阻越大，理论阻尼系数 \(\beta=R_\Sigma/(2L)\) 越大，振幅衰减越快。本文拟合值 \(9.54\times10^3\ \mathrm{s^{-1}}\) 高于主写理论值 \(7.25\times10^3\ \mathrm{s^{-1}}\)，说明实际阻尼可能大于 \(150\ \Omega+140\ \Omega\) 的估计。
    \item 电感直流内阻和交流损耗不完全相同。万用表测得的 \(R_L=92.825\ \Omega\) 是直流内阻，而实验频率约为 \(11\ \mathrm{kHz}\)。在交流条件下，线圈损耗、磁芯损耗和连接接触电阻可能使等效损耗大于直流读数。
    \item 峰谷读数和周期判断会显著影响阻尼系数。若把相邻记录点误作半个周期，会使 \(\beta\) 被放大约 2 倍。低电压尾端点还容易受示波器分辨率和噪声影响；若末端峰值读数偏小，\(\ln(V_C)\) 会偏低，拟合斜率绝对值变大，导致 \(\beta_{\mathrm{fit}}\) 偏大。
    \item 半功率点由有限采样和曲线读数得到。\(f_+=15.4\ \mathrm{kHz}\) 不在频响表中，说明它来自曲线读数或表外记录。若 \(f_+\) 被读得偏高或 \(f_-\) 被读得偏低，则 \(\Delta f\) 偏大，\(Q=f_0/\Delta f\) 偏小；反之 \(Q\) 偏大。
    \item 信号源内阻会影响分压估算。若 \(V_m\) 不是负载端实际电压，而包含信号源内阻压降或开路设定值，则 \(V_m/V_R\) 会偏离真实分压比，进而影响 \(R_{\mathrm{loss},f}\) 的反推值和功率账。
    \item 元件参数公差会影响理论共振频率。理论值 \(f_{0,\mathrm{th}}\) 使用标称 \(L\) 和 \(C\)。实验值 \(11.4\ \mathrm{kHz}\) 比理论值 \(11.254\ \mathrm{kHz}\) 高约 \(1.30\%\)，该偏差处于元件参数、曲线读数和仪器分辨率共同作用的范围内。
\end{enumerate}

\section{实验小结}

本实验观察了 RLC 串联电路的阻尼振荡和受迫振荡特性。欠阻尼部分中，逐次峰谷电压的对数与测量序号呈良好线性关系，拟合式为
\[
    \ln(V_C)=-0.8132n+2.3187,\qquad r=-0.9998.
\]
将相邻记录点间隔取为一个周期 \(T=85.2\ \mu\mathrm{s}\)，得到
\[
    \beta_{\mathrm{fit}}=9.54\times10^3\ \mathrm{s^{-1}},\qquad
    \tau=104.8\ \mu\mathrm{s}.
\]
用阻尼振荡原始损耗 \(R_{\mathrm{loss},d}=140\ \Omega\) 和外接电阻 \(150\ \Omega\) 计算，理论阻尼系数为
\[
    \beta_{\mathrm{th}}=7.25\times10^3\ \mathrm{s^{-1}}.
\]
实验值比理论值高约 \(31.6\%\)，说明实际损耗和读数误差会使衰减快于简化模型预测。临界阻尼理论电阻为
\[
    R_c=2.83\ \mathrm{k}\Omega.
\]

受迫振荡部分中，幅频响应曲线在 \(11.4\ \mathrm{kHz}\) 附近达到峰值，相频曲线在同一区域穿过零点。理论共振频率为
\[
    f_{0,\mathrm{th}}=11.254\ \mathrm{kHz},
\]
实验值的相对误差为 \(1.30\%\)。由半功率频率 \(f_-=8.7\ \mathrm{kHz}\)、\(f_+=15.4\ \mathrm{kHz}\) 得到
\[
    \Delta f=6.7\ \mathrm{kHz},\qquad Q=1.70.
\]
其中 \(f_+=15.4\ \mathrm{kHz}\) 来自曲线读数或表外记录，而非表中直接读数。

损耗电阻处理是本实验数据分析中的关键点。\(140\ \Omega\) 是阻尼振荡原始损耗 \(R_{\mathrm{loss},d}\)，\(194.3\ \Omega\) 是按教师批注 \(R_m=500\ \Omega\) 由分压公式得到的受迫振荡估算损耗 \(R_{\mathrm{loss},f}\)。两者来源不同，不能合并成同一个寄生电阻。

功率核算统一采用峰值体系后，得到
\[
    I_m=6.656\ \mathrm{mA},\quad
    P_R=11.08\ \mathrm{mW},\quad
    P_{\mathrm{loss},f}=4.30\ \mathrm{mW},\quad
    P_{\mathrm{source}}=15.38\ \mathrm{mW}.
\]
使用 \(R_{\mathrm{loss},f}\) 时，\(P_R+P_{\mathrm{loss},f}\) 与 \(P_{\mathrm{source}}\) 构成功率账自洽；使用独立记录 \(R_{\mathrm{loss},d}=140\ \Omega\) 时，功率账相差约 \(7.8\%\)。因此，分压估算损耗带来的功率闭合应理解为功率账自洽，而不是完全独立的检验。

\section{思考题}

\subsection*{1. 证明 \(\sqrt{L/C}\) 的量纲与电阻相同}

由电感定义 \(V=L\,\mathrm{d}I/\mathrm{d}t\) 可知
\[
    [L]=\frac{\mathrm{V}\cdot\mathrm{s}}{\mathrm{A}}=\Omega\cdot\mathrm{s}.
\]
由电容定义 \(Q=CV\) 可知
\[
    [C]=\frac{\mathrm{A}\cdot\mathrm{s}}{\mathrm{V}}=\frac{\mathrm{s}}{\Omega}.
\]
因此
\[
    \left[\sqrt{\frac{L}{C}}\right]
    =\sqrt{\frac{\Omega\cdot\mathrm{s}}{\mathrm{s}/\Omega}}
    =\sqrt{\Omega^2}
    =\Omega.
\]
所以 \(\sqrt{L/C}\) 与电阻具有相同量纲。

\subsection*{2. 信号发生器接入 RLC 电路后方波变形的原因}

方波可以分解为基波和一系列奇次谐波。信号发生器可等效为理想电压源串联内阻 \(R_s\)，通常约为 \(50\ \Omega\)。未接负载或负载阻抗很高时，\(R_s\) 上压降很小，示波器显示的方波较完整。接入 RLC 电路后，负载阻抗
\[
    Z(\omega)=R+j\left(\omega L-\frac{1}{\omega C}\right)
\]
随频率变化，不同谐波分量在 \(R_s\) 与 \(Z(\omega)\) 组成的分压器中衰减比例不同。低频和高频谐波受电容、电感阻抗影响较大，接近共振频率的分量传输较强，因此输出波形不再保持理想方波，而会出现畸变、振铃或边沿过渡变化。

\subsection*{3. 比较不同电阻值下频率响应曲线的差异}

在相同 \(L\)、\(C\) 下，串联 RLC 电路的品质因数近似满足
\[
    Q=\frac{\omega_0L}{R_\Sigma},
    \qquad
    \Delta\omega=\frac{R_\Sigma}{L}.
\]
因此，总电阻越大，带宽越宽，峰值越低，频率选择性越差。若忽略额外损耗、仅作教学讨论，取 \(\omega_0L\approx1414\ \Omega\)，则表~\ref{tab:question-q} 给出不同电阻下的近似变化趋势。

\begin{table}[H]
\centering
\caption{不同电阻下频率响应曲线的近似变化}
\label{tab:question-q}
\begin{tabularx}{\textwidth}{>{\raggedleft\arraybackslash}p{3cm} >{\centering\arraybackslash}p{2.5cm} Y}
\toprule
电阻 & 近似 \(Q\) & 曲线特征 \\
\midrule
\(1.25\ \mathrm{k}\Omega\) & 1.13 & 峰较尖，频率选择性较好，电容电压可能出现较明显谐振放大 \\
\(2.5\ \mathrm{k}\Omega\) & 0.57 & 峰较平，选择性减弱 \\
\(4.0\ \mathrm{k}\Omega\) & 0.35 & 带宽更大，峰值更低，共振特征不明显 \\
\bottomrule
\end{tabularx}
\end{table}

以上讨论的是不同电阻条件下频率响应变化的一般趋势，与前文原始记录中 \(R_f=2.5\ \mathrm{k}\Omega\) 和 \(R_m=500\ \Omega\) 的分工不混写。

\subsection*{4. 改变步骤 4.2 中 6、7 的电阻值时，频率响应曲线如何变化}

若步骤 6 对应带阻或陷波滤波，减小电阻会提高回路的选择性，使陷波更窄、更深；增大电阻会使陷波变宽变浅，滤波效果变差。若步骤 7 对应并联或谐振响应，减小电阻性损耗会提高品质因数，使共振峰更尖锐、峰值更高；增大电阻性损耗会降低品质因数，使峰值下降、带宽增大。两种情况的本质都是电阻通过改变 \(Q\) 来改变频率选择性。

\subsection*{5. 为什么本实验通常不需要长时间等待“振荡稳定后”再记录数据}

RLC 电路的暂态建立时间由电路时间常数决定，通常远小于人工观察和示波器刷新时间。在本实验量级下，若按受迫振荡中的总电阻取 \(500\sim2500\ \Omega\)，暂态时间常数
\[
    \tau\sim\frac{2L}{R_\Sigma}
\]
约为几十微秒量级，即使按 \(5\tau\) 估计，建立时间也通常小于 \(0.5\ \mathrm{ms}\)。因此本实验中示波器上看到的受迫响应通常已经接近稳态，不必长时间等待；但在精密测量时，仍应保证波形触发稳定后再读数。

\section{参考文献}

\begin{enumerate}[label={[\arabic*]}, leftmargin=2.4em, itemsep=0.4em]
    \item 清华大学物理系. 基础物理实验讲义：RLC 电路[Z]. 北京：清华大学, 2024.
    \item Karni S. \textit{Applied Circuit Analysis} [M]. New York: John Wiley \& Sons Inc., 1988: 220--233.
\end{enumerate}

\appendix
\section{原始数据记录摘录}

本附录整理实验记录中的关键原始数据，便于查证正文中的计算过程。正文中涉及的拟合、共振频率、品质因数和功率核算均由以下数据及相应公式得到。

\begin{table}[H]
\centering
\caption{电源操作原始记录}
\begin{tabular}{cccc}
\toprule
电阻 \(R/\Omega\) & 电压 \(V/\mathrm{V}\) & 电流 \(I/\mathrm{A}\) & 工作状态 \\
\midrule
30 & 3.19 & 0.100 & CC，恒流 \\
40 & 4.18 & 0.100 & CC，恒流 \\
50 & 5.00 & 0.095 & CV，恒压 \\
60 & 5.00 & 0.081 & CV，恒压 \\
70 & 5.00 & 0.070 & CV，恒压 \\
\bottomrule
\end{tabular}
\end{table}

\begin{table}[H]
\centering
\caption{阻尼振荡峰谷电压原始记录与对数值}
\begin{tabular}{cccc}
\toprule
测量次数 \(n\) & \(V_C/\mathrm{V}\) & \(\ln(V_C)\) & 时间记录 \\
\midrule
1 & 4.7180 & 1.5514 & \(76\ \mu\mathrm{s}\) \\
2 & 1.9532 & 0.6695 & -- \\
3 & 0.8437 & \(-0.1700\) & -- \\
4 & 0.3906 & \(-0.9401\) & -- \\
5 & 0.1781 & \(-1.7254\) & -- \\
6 & 0.0781 & \(-2.5498\) & \(502\ \mu\mathrm{s}\) \\
\bottomrule
\end{tabular}
\end{table}

\begin{table}[H]
\centering
\caption{受迫振荡频率响应原始记录}
\begin{tabular}{ccc}
\toprule
频率 \(f/\mathrm{kHz}\) & 电阻电压峰值 \(V_{Rm}/\mathrm{V}\) & 相位差 \(\varphi/^\circ\) \\
\midrule
5.0 & 0.933 & \(-73.5\) \\
6.0 & 1.211 & \(-69.9\) \\
7.0 & 1.600 & \(-64.4\) \\
8.0 & 2.029 & \(-55.8\) \\
9.0 & 2.641 & \(-43.2\) \\
10.0 & 3.017 & \(-26.1\) \\
10.5 & 3.187 & \(-11.6\) \\
11.0 & 3.287 & \(-3.7\) \\
11.5 & 3.311 & 4.0 \\
12.0 & 3.268 & 9.4 \\
13.0 & 3.055 & 22.9 \\
14.0 & 2.770 & 17.2 \\
15.0 & 2.502 & 23.6 \\
\bottomrule
\end{tabular}
\end{table}

\begin{table}[H]
\centering
\caption{共振点及半功率点相关记录}
\begin{tabularx}{\textwidth}{>{\centering\arraybackslash}p{4cm} C Y}
\toprule
项目 & 记录值 & 说明 \\
\midrule
实验共振频率 \(f_0\) & \(11.4\ \mathrm{kHz}\) & 由曲线读数或插值得到 \\
输入电压峰值 \(V_m\) & \(4.621\ \mathrm{V}\) & 共振附近记录 \\
电阻电压峰值 \(V_R\) & \(3.328\ \mathrm{V}\) & 用于电流和功率核算 \\
半功率低频点 \(f_-\) & \(8.7\ \mathrm{kHz}\) & 原始记录/曲线读数 \\
半功率高频点 \(f_+\) & \(15.4\ \mathrm{kHz}\) & 表外原始记录或曲线读数 \\
阻尼振荡损耗 \(R_{\mathrm{loss},d}\) & \(140\ \Omega\) & 用于阻尼理论比较 \\
受迫估算损耗 \(R_{\mathrm{loss},f}\) & \(194.3\ \Omega\) & 由 \(R_m=500\ \Omega\) 与 \(V_m/V_R\) 反推 \\
\bottomrule
\end{tabularx}
\end{table}

\end{document}

